\section{Data Access Services} \label{sec:services}

\subsection{Community Alert Brokers} \label{sec:brokers}

Near real-time access to alert packets is available through services provided by community alert brokers.
Due to limited bandwidth from the USDF, seven brokers were selected through a proposal process \citep{LDM-612} to receive the alert stream directly from the alert distribution system at the USDF (\S \ref{sec:alertdist}):  

\begin{itemize}
    \item ALeRCE: Automatic Learning for the Rapid Classification of Events \citep{Forster:21:ALeRCE}
    \item AMPEL: Alert Management, Photometry, and Evaluation of Light Curves \citep{Nordin:19:AMPEL}
    \item ANTARES: Arizona-NOIRLab Temporal Analysis and Response to Events System \citep{Matheson:21:ANTARES}
    \item Babamul  \citep{2025arXiv251100164J}
    \item Fink \citep{Moller:21:FINK}
    \item Lasair \citep{Smith:19:Lasair}
    \item Pitt-Google
\end{itemize}

As alert packets are world-public and freely sharable (\S \ref{sec:data-rights}), additional brokers and other services may operate ``downstream'' by connecting to full or filtered streams provided by the upstream brokers.
As of this writing, downstream brokers include

\begin{itemize}
    \item Point of Interest Broker
    \item SNAPS: Solar System Notification Alert Processing System 
\end{itemize}

\eric{check broker sites if there are newer citations}

\subsection{Alert Filtering Service}

The original requirements on the Rubin Data Management system (DMS) specify that a ``basic, limited capacity, alert filtering service shall be provided that can be given user defined filters to reduce the alert stream to manageable levels" and that users of this service ``shall be able to use a predefined set of simple filters" \citep{LSE-61}.
These requirements were written before the current brokers existed, and they are already providing beyond this basic required functionality.
Instead of creating a redundant alert filtering service, Rubin staff are instead supporting the community to use alerts by creating and mainting a set of ``community filters", to be installed in the NOIRLab's ANTARES broker (chosen because of the NOIRLab connection).
As described in \citet{RTN-090}, these community filters will cover a range of basic science use-cases and be generated based on community input.
The first community filters are expected to be installed and validated by the fall of 2026.

\subsection{Rubin Science Platform}

\gregory{todo}

\subsubsection{Alert Archive}

Each alert packet is stored as-transmitted in the USDF.
\citet{DMTN-183} describes the motivation for and implementation of the archive.
An ingester service listens to the Kafka alert stream and writes each packet to a key-value store indexed by the alert ID.
A front-end service, \texttt{herald} \citep{SQR-114}, allows authenticated users to retrieve individual alerts in Avro Open Container Format, JSON, or FITS.
The Portal aspect of the Rubin Science Platform provides an Alert Visualizer which displays the image cutouts and catalog data.

\subsubsection{Image Services}

\subsubsection{Butler Data Products}

\subsubsection{Prompt Products Database} \label{sec:ppdb}

\subsection{IDACs}

\todo{confirm IDAC section}